% Desenvolvido por André Leite - DE - UFPE
% Ajustes Raydonal Ospina - DE - UFPE/UFBA

\documentclass[a4paper,11pt,oneside,twocolumn]{Config/milktest}

%---------------------- customização ----------------------



\usepackage{fontspec}
\usepackage{xunicode}
\usepackage{xltxtra}
\usepackage[LGR,T1]{fontenc}
%\usepackage{ebgaramond}

%\usepackage[cp1252]{inputenc}                %% Use um dos dois esquemas
\usepackage[english,brazil]{babel}                 %% Hifenizacao e estilo
\usepackage{pstricks}                          %% Pacote de comandos p/ Facilitar o Postscripts
\usepackage{pstricks-add}                    %% " "
\usepackage{pst-all,pst-math,pst-plot, pst-func}       %% " "
\usepackage{textcomp}							%%
%\usepackage[tight,nice]{units}				%%
\usepackage[centertags]{amsmath}       %% [leqno]
\usepackage{amsfonts}                         %% Pacotes da AMS
\usepackage{amssymb}                         %%
\usepackage{stmaryrd} 							%%
\usepackage{Config/milkstyle}
%\usepackage{sabon}

\usepackage{nicefrac}%  Nice Frescura
\usepackage{balance}%
\usepackage{lastpage} %                      %% Referênciar a última página
\usepackage{graphicx}                        %% Mais ferramentas para
%\usepackage[T1]{pbsi} % Procurar e instalar no Mac!!!!! <------<-------<------
\usepackage{mathrsfs}                        %% $\mathscr{ABCD}$
\usepackage{lettrine}                        %% ...
%\usepackage[normal,bf,up]{caption2}          %% Controlar melhor os captions....
\usepackage{eucal}                           %% beau
\usepackage{colortbl} %
\usepackage{hyperref}                 %%
%\usepackage{amsthm}
\usepackage{indentfirst}%
\usepackage{enumitem}%
\usepackage{colordvi}                        %% Usar Cores...
\usepackage{fancyhdr}%
\usepackage{yfonts}%
\usepackage{theorem}%
\usepackage{multirow}                        %% Juntar linhas em tabelas
\usepackage{midfloat}%
\usepackage{cancel}%
%\usepackage[square]{natbib}%sectionbib,
\usepackage{fancyhdr} % cabeçalhos ...
\usepackage{rotating}


\usepackage{lipsum} 
%% Lista de Pacotes utilizados


%% Deixando o caption mais interessante...
%\renewcommand{\captionfont}{\itshape}


\hypersetup{%backref,%
%            pdfpagemode=FullScreen,%
            colorlinks=true,%
            citecolor=blue,%
            urlcolor=blue, %
            % urlcolor=black, %
            linkcolor=blue,%
            pdfauthor={André Leite Wanderley}, %
            pdftitle={1EE/ET406}, %
            pdfcreator={André Leite Wanderley}, %
            breaklinks=true, %
            linktocpage=true
            }%


\DeclareMathOperator*{\Inf}{Inf}%
\DeclareMathOperator*{\Sup}{Sup}%
\DeclareMathOperator*{\Min}{Min}%
\DeclareMathOperator*{\Max}{Max}
\DeclareMathOperator*{\Sgn}{Sgn}%
\DeclareMathOperator*{\E}{E}%
\newcommand{\mbf}{\boldsymbol}%
\DeclareMathAlphabet\mathbb{U}{fplmbb}{m}{n}%
%
%\numberwithin{equation}{section}                       %%
%\numberwithin{table}{section}                          %% Numeração dos floats...
%\numberwithin{figure}{section}                         %%


\newcommand{\oo}{\mathbin{\bullet}}
\newcommand{\oc}{\mathbin{\circ}}
\newcommand{\noo}{\mathbin{\not\!\bullet}}
\newcommand{\cf}{\emph{cf.}\ }
\newcommand{\ie}{\emph{i.e.}}
\newcommand{\eg}{\emph{e.g.}}
\newcommand{\raab}{\emph{reductio ad absurdum}}
\newcommand{\eetc}{\textit{et c\ae tera}}
\newcommand{\dom}{\mathfrak{D}}
\newcommand{\ima}{\mathfrak{I}}
\newcommand{\idenA}{\mathcal{I}_\mathcal{A}}
\newcommand{\ce}{\ensuremath{\mathcal{X}}}
\newcommand{\piq}{\ensuremath{\succsim}}
\newcommand{\pq}{\ensuremath{\succ}}
\newcommand{\iq}{\ensuremath{\sim}}
\newcommand{\niq}{\ensuremath{\nsim}}
\newcommand{\npiq}{\ensuremath{\not\succsim}}
\newcommand{\npq}{\ensuremath{\nsucc}}
\newcommand{\po}[2]{\ensuremath{(#1,#2)}}

%\setlist{noitemsep}
%\setlist{labelindent=\parindent} % << Usually a good idea
\setlist{itemsep=0pt}%
\setitemize{leftmargin=\parindent} %
\setitemize[1]{label=$\triangleright$} %
%\setenumerate[1]{label=\arabic*.,ref=\arabic*}
\setenumerate[1]{label=\arabic*.,ref=\arabic*}
\setenumerate[2]{label=\theenumi.\arabic*.,ref=\theenumi.\arabic*}
\setenumerate[3]{label=\roman*.,ref=\theenumii.\roman*}
%\setdescription{%font=\sffamily\bfseries}

%                }%


%%%%%%%%%%%%%%%%%%%%%%%%%%%%%%%%
% Para facilitar minha vida

\newcommand{\Ds}{\displaystyle}%
\newcommand{\Ts}{\textstyle}%
\newcommand{\Ss}{\scriptstyle}%
\newcommand{\SSs}{\scriptscriptstyle}%

\def\aureos{%
            \newrgbcolor{gray1}{.9 .9 .9}
            \newrgbcolor{guess1}{1 0.388235294 0.090196078}
            \psset{unit=1in,cornersize=absolute}%
            \psline[linecolor=guess1,linewidth=1.2pt](-0.125,-1.291667)(1.1125,-1.291667)
            \psline[linecolor=guess1,linewidth=0.6pt](0.347683,-1.291667)(0.347683,-0.52685)
            \psline[linecolor=guess1,linewidth=0.6pt](0.347683,-0.999533)(-0.125,-0.999533)
            \psline[linecolor=guess1,linewidth=0.6pt](0.167134,-0.999533)(0.167134,-1.291667)
            \psline[linecolor=guess1,linewidth=0.6pt](0.167134,-1.111118)(0.347683,-1.111118)
            \psline[linecolor=guess1,linewidth=0.6pt](0.236098,-1.111118)(0.236098,-0.999533)
            \psline[linecolor=guess1,linewidth=0.6pt](0.236098,-1.068496)(0.167134,-1.068496)
            \psline[linecolor=guess1,linewidth=0.6pt](0.209756,-1.068496)(0.209756,-1.111118)
            \psline[linecolor=guess1,linewidth=0.6pt](0.209756,-1.084776)(0.236098,-1.084776)
           }%

\def\GuEsS{%
            \newrgbcolor{guess1}{1 0.388235294 0.090196078}
            \newrgbcolor{guess2}{0.960784314 0.109803922 0.043137255}
            \psset{unit=1in,cornersize=absolute}%
            \rput(0,-1.5){{\fontsize{35}{0}\usefont{OT1}{cmr}{m}{n}\selectfont\textsc{\textbf{S}}}}%
            \rput(0.625,-1.5){\shortstack{{\huge\textbf{\textsc{ystems}}}\\
                {\scriptsize\sc\hfill Research Group}}}
            \psline[linecolor=black,fillcolor=guess2,fillstyle=none,linewidth=3.6pt](-0.125,-1.725)(1.1125,-1.725)%
            \uput{0.5ex}[l](1.148214,-1.229167){\llap{{\tiny\sc UFBA}}}%
            \psline[linecolor=guess1,fillcolor=guess1,fillstyle=none,linewidth=1.2pt](-0.125,-1.291667)(1.1125,-1.291667)%
            \psline[linecolor=guess1,fillcolor=guess1,fillstyle=none,linewidth=.6pt](0.347683,-1.291667)(0.347683,-0.52685)%
            \psline[linecolor=guess1,fillcolor=guess1,fillstyle=none,linewidth=.6pt](0.347683,-0.999533)(-0.125,-0.999533)%
            \psline[linecolor=guess1,fillcolor=guess1,fillstyle=none,linewidth=.6pt](0.167134,-0.999533)(0.167134,-1.291667)%
            \psline[linecolor=guess1,fillcolor=guess1,fillstyle=none,linewidth=.6pt](0.167134,-1.111118)(0.347683,-1.111118)%
            \psline[linecolor=guess1,fillcolor=guess1,fillstyle=none,linewidth=.6pt](0.236098,-1.111118)(0.236098,-0.999533)%
            \psline[linecolor=guess1,fillcolor=guess1,fillstyle=none,linewidth=.6pt](0.236098,-1.068496)(0.167134,-1.068496)%
            \psline[linecolor=guess1,fillcolor=guess1,fillstyle=none,linewidth=.6pt](0.209756,-1.068496)(0.209756,-1.111118)%
            \psline[linecolor=guess1,fillcolor=guess1,fillstyle=none,linewidth=.6pt](0.209756,-1.084776)(0.236098,-1.084776)%
           }%

%%%%%%%%%%%%%%%%%%%%%%%%%%%%%%%%%%%%%%%%%%%%%%%%%%%%%%%%


\theoremstyle{plain}                     %%
%\theoremstyle{margin}                   %%
\theoremheaderfont{\scshape\bfseries}    %%
\theorembodyfont{\rmfamily}                       %%

\newtheorem{Prop}{\small Proposition}[section]  %%
\newtheorem{Def}{\small Definition}[section]    %%
\newtheorem{Not}{\small Notice}[section]    %%
\newtheorem{Theo}{\small Theorem}[section]      %%
\newtheorem{Lem}{\small Lema}[section]         %%
\newtheorem{Coro}{\small Corollary}[section]    %%
\newtheorem{Conj}{\small Conjecture}[section]   %%
\newtheorem{Prob}{\small Problem}[section]      %%
\newtheorem{Exa}{\small Example}[section]       %%
%\theoremheaderfont{\bfseries}    %%
%\theorembodyfont{\rmfamily}  %%
\newtheorem{proof}{\small Demonstration}[section]
\newtheorem{sol}{\small Solution}[section]


\newenvironment{Proofm}               % Deve ter uma maneira
    {%\strut \par\medskip             % ridiculamente mais
     %\strut\hspace{\parindent}      % fácil de fazer isso...
     %\begin{minipage}{.8\linewidth} % I'm sure about that.
     \begin{proof}\small}                   %
    {\strut\hfill\textsc{q.e.d.}\end{proof}                      %
    %\end{minipage}                  %
    %\strut\par\medskip
                        }             %

\newenvironment{solm}               % Deve ter uma maneira
    {%\strut\par%\medskip             % ridiculamente mais
     %\strut\hspace{\parindent}      % fácil de fazer isso...
     %\begin{minipage}{.8\linewidth} % I'm sure about that.
     %\begin{sol}
     \small
     \noindent\textbf{Solução:}
     %\par %\smallskip
     }                   %
    {\strut\hfill\textsc{$\blacksquare$}%\end{sol}                      %
    %\end{minipage}                  %
    %\strut\par\medskip
                        }             %



\newenvironment{Propm}               % Deve ter uma maneira
    {\strut\par\smallskip             % ridiculamente mais
     \strut\hspace{\parindent}      % fácil de fazer isso...
     \begin{minipage}{.8\linewidth} % I'm sure about that.
     \begin{Prop}\small}                   %
    {\end{Prop}                      %
    \end{minipage}                  %
    \strut\par\medskip
    }             %


\newenvironment{Defm}               % Deve ter uma maneira
    {\strut\par\smallskip             % ridiculamente mais
     \strut\hspace{\parindent}      % fácil de fazer isso...
     \begin{minipage}{.8\linewidth} % I'm sure about that.
     \begin{Def}\small}                   %
    {\end{Def}                      %
    \end{minipage}                  %
    \strut\par\medskip}             %

\newenvironment{Theom}               % Deve ter uma maneira
    {\strut\par\smallskip             % ridiculamente mais
     \strut\hspace{\parindent}      % fácil de fazer isso...
     \begin{minipage}{.8\linewidth} % I'm sure about that.
     \begin{Theo}\small}                   %
    {\end{Theo}                      %
    \end{minipage}                  %
    \strut\par\medskip}

\newenvironment{Notm}               % Deve ter uma maneira
    {\strut\par\smallskip             % ridiculamente mais
     \strut\hspace{\parindent}      % fácil de fazer isso...
     \begin{minipage}{.8\linewidth} % I'm sure about that.
     \begin{Not}\small}                   %
    {\end{Not}                      %
    \end{minipage}                  %
    \strut\par\medskip}


\newenvironment{Exam}               % Deve ter uma maneira
    {%\strut\par\smallskip             % ridiculamente mais
     %\strut\hspace{\parindent}      % fácil de fazer isso...
     %\begin{minipage}{1\linewidth} % I'm sure about that.
     \begin{Exa}\small}                   %
    { \scriptsize$\blacksquare$\end{Exa}                      %
    %\end{minipage}                  %
    \strut\par\medskip}

\def\Pt#1(#2,#3){
       \expandafter\def\csname #1\endcsname{#2,#3}%
       \expandafter\def\csname #1X\endcsname{#2}%
       \expandafter\def\csname #1Y\endcsname{#3}}%

\newcommand{\hlinem}[3]{\noalign{\vspace{#3}\hrule height #1 depth #2}}
\newcommand{\hlinemilk}[3]{\noalign{}}

\newcommand{\quem}[1]{\par\medskip\strut\hfill\bf  --- #1}
\newcounter{questionNumber}
\newcommand{\question}{\stepcounter{questionNumber}\lettrine[lhang=.75, findent=10pt, nindent=0pt,lines=3, lraise=.2]{\thequestionNumber}{$\blacktriangleright$}}


\newenvironment{milkquote}
    {\medskip %
    \begin{center}
    \begin{minipage}{.9\linewidth}
    \hspace{.25cm}\small\itshape}%
    {\medskip\end{minipage}
    \end{center}
    }%


\newcommand{\fim}{%
                 \par\medskip
                 \begin{center}
                      \includegraphics[width=1cm]{ornaments/ornamento4.eps}
                 \end{center}
                 \vspace*{.5cm}
                 }%



\newenvironment{milkanswer}
    {\par %
    \begin{center}
    \color{orange}\textbullet\quad\textbf{Solution}\quad\textbullet
    \end{center}\small
    }%
    {
    \fim
    }


%%%
%%% Set up page layout parameters.
%%%
\setlength\paperheight{297mm}      %
\setlength\paperwidth{210mm}       % papel a4
%\setlength{\parindent}{6mm}

\addtolength{\hoffset}{-1in}%
\addtolength{\voffset}{-1.5in}

\setlength{\oddsidemargin}{25mm}    % 
\setlength{\topmargin}{5cm}        %
\setlength{\textwidth}{170mm}      % Margem direita: 25mm
\setlength{\textheight}{220mm}  % Margem Inferior: 25mm
\setlength{\marginparsep}{0mm}     %

\setlength{\headheight}{16mm}
\renewcommand{\columnsep}{1cm}%
%\marginsize{1.25in}{1.75in}{1.25in}{1.25in}



% If you can to leave the header here
%---

%---

\renewcommand{\headrulewidth}{0pt}
\definecolor{cor}{rgb}{.5,0, 0}

%%%%%%%%%%%%%%%%%%%%%%%%%%% Font 

\defaultfontfeatures{Ligatures=Historic,Contextuals=Alternate, Numbers=OldStyle,RawFeature={ +hlig, +dlig, +ss05}}
\setmainfont[RawFeature={+ss05, +hlig, +dlig}, BoldFont = Gentium Basic Bold, ItalicFeatures={RawFeature={+clig},CharacterVariant=5:0}]{EB Garamond}
% % Do the replacements manually in titles, to not ruin the textsc
% \newfontfamily\booktitlefont[RawFeature={-ss02},LetterSpace=40,WordSpace=6]{EB Garamond}
% \newfontfamily\spacedfont[RawFeature={-ss02},LetterSpace=20,WordSpace=3]{EB Garamond}
% \newfontfamily\lettrinefont{EB Garamond Initials}
% \newfontfamily\headerfont{EB Garamond}

%\setmainfont[BoldFont = Gentium Basic Bold,
%      ItalicFont = Gentium Plus Italic,
%      BoldItalicFont = Gentium Basic Bold Italic,
%      Mapping = tex-text]{Gentium Plus}





 

\usepackage{bm}
\usepackage{verbatimbox}
\usepackage{float}

%%%%%%%%%%%%%----- Header------%%%%%%%%%%%%%%%%%%%
\newsavebox{\mygraphic}
\sbox{\mygraphic}{\includegraphics[scale=1.2]{Figuras/ufba.eps}}

\fancypagestyle{milktablestyle}{
    \fancyhead{}
    \fancyfoot[C]{\thepage/\pageref{LastPage}} %
}

\pagestyle{fancy}

\fancyhead{}
\fancyhead[C]{
\rput(0,2){
  \begin{minipage}{1.75cm}\centering
    \includegraphics[width=1.75cm]{Figuras/ufbaBlack.eps}
  \end{minipage}~\hfill~\begin{minipage}{10cm}\Large\centering
   \color{cor}{\scshape universidade federal da bahia}\\
    {\scshape instituto de matemática e estatística}\\
    {\scshape departamento de estatística}\\[5pt]
    % Aqui as informações da prova e a data
    {\Large\scshape GABARITO -  prova I --- matd44 --- 17/06/2025} 
  \end{minipage}~\hfill~\begin{minipage}{1.75cm}\centering
    \includegraphics[width=1.75cm]{Figuras/ufba.eps}
  \end{minipage}
  }}

\fancyfoot[C]{\thepage/\pageref{LastPage}} % 

%%%%%%%%%%%%%----- Header------%%%%%%%%%%%%%%%%%%%

% \usepackage[printwatermark]{xwatermark}
% \newwatermark[allpages,color=gray!20,angle=45,scale=3,xpos=30,ypos=-60]{DRAFT}

\begin{document}




% %------------------------------------------------
\beb
{\small
\begin{description}
\item[Professor:] Raydonal Ospina
% \item[E-mail:] raydonalmartinez@ufba.br
% \item[Site do curso:] \url{https://castlab.org/courses/matd49/}

\item[Regras:] A prova é individual. Leia com atenção as perguntas. A prova deve ser claramente resolvida. 
% Seja claro e organizado. Pode fazer uso de computadora, calculadora, celular, I.A., livros, material de aula.

% \item[Atenção:] Responda única e exclusivamente o número de questões {\bf suficientes} para completar 10 (dez) pontos.  Fique a vontade para escolher as perguntas que achar mais interessantes. Especifique sua escolha, por exemplo, 1(a). (isto quer dizer que você escolheu  o numeral (a) da questão 1).
\end{description}
}
\eeb

% %%%%%%%%%%%%%%%%%%%%%%%%%%%%%%%%%%%%%%%%%%%%%%%%%%%%%%%%%%%%%%%%%
\balance
%%%%%%%%%%%%%%%%%%%%%%%%%%%%%%%%%%%%%%%%%%%%%%%%%%%%%%%%%%%%%%%%%




\medskip 
\question~({\bf Plano amostral geral})    Seja 
${\cal U}=\{1,2,3\}$ uma população finita de tamanho $N=3$ e $\mathbf{Y} =\{1,2,3\}$  o vetor da característica populacional renda bruta (mensal em salários mínimos) familiar.  Suponha que o seguinte plano amostral é implementado 
$p(s_1)=p(\{1,2\})=\frac{1}{2},$ $p(s_2)=p(\{1,3\})=\frac{1}{4}$ e
$p(s_3)=p(\{2,3\})=\frac{1}{4}.$
\begin{description}
% \item[a)]~\textbf{[2pts]} Determine as probabilidades de inclusão de primeira ordem. 
\item[a)] Determine as probabilidades de inclusão de primeira e  segunda ordem. Determine se o plano amostral induzido pelo esquema de amostragem proposto é mensurável?


{\scriptsize
\noindent \textbf{Res: 1a.} Primeiro, vamos definir as probabilidades de inclusão de primeira ordem. A probabilidade de inclusão de primeira ordem \(\pi_i\) é a probabilidade de um elemento \(i\) ser incluído na amostra. 

Para \(i = 1\):
\[
\pi_1 = P(\{1, 2\}) + P(\{1, 3\}) = \frac{1}{2} + \frac{1}{4} = \frac{3}{4}
\]

Para \(i = 2\):
\[
\pi_2 = P(\{1, 2\}) + P(\{2, 3\}) = \frac{1}{2} + \frac{1}{4} = \frac{3}{4}
\]

Para \(i = 3\):
\[
\pi_3 = P(\{1, 3\}) + P(\{2, 3\}) = \frac{1}{4} + \frac{1}{4} = \frac{1}{2}
\] 
A probabilidade de inclusão de segunda ordem \(\pi_{ij}\) é a probabilidade de os elementos \(i\) e \(j\) serem incluídos simultaneamente na amostra.

Para \(\pi_{12}\):
\[
\pi_{12} = P(\{1, 2\}) = \frac{1}{2}
\]

Para \(\pi_{13}\):
\[
\pi_{13} = P(\{1, 3\}) = \frac{1}{4}
\]

Para \(\pi_{23}\):
\[
\pi_{23} = P(\{2, 3\}) = \frac{1}{4}
\]

Um plano amostral é mensurável se todas as probabilidades de inclusão de segunda ordem forem maiores que zero e se \(\pi_{ij} \leq \min(\pi_i, \pi_j)\). Neste caso, temos:

\[
\pi_{12} = \frac{1}{2} \quad \text{e} \quad \min(\pi_1, \pi_2) = \frac{3}{4}
\]
\[
\pi_{13} = \frac{1}{4} \quad \text{e} \quad \min(\pi_1, \pi_3) = \frac{1}{2}
\]
\[
\pi_{23} = \frac{1}{4} \quad \text{e} \quad \min(\pi_2, \pi_3) = \frac{1}{2}
\]

Como todas as probabilidades de inclusão de segunda ordem são menores ou iguais ao mínimo das probabilidades de inclusão de primeira ordem correspondentes, podemos concluir que o plano amostral é mensurável.
}


% \item[c)]~\textbf{[2pts]} Determine as covariâncias dos indicadores de pertencimento à amostra  $$I_k=I_k(s)=\begin{cases} 1, & k \in s, \\ 0, & \text{caso contrário}.\end{cases}$$
\item[b)] Forneça a distribuição de probabilidades do estimador de Horvitz-Thompson $\hat{t}_\pi$ para o total populacional da renda bruta (mensal em salários mínimos) familiar.

\medskip
{\scriptsize

\textbf{Res:  1b} O estimador de Horvitz-Thompson para o total populacional da renda bruta (mensal em salários mínimos) familiar é dado por:

\[
\hat{t}_\pi = \sum_{i \in s} \frac{y_i}{\pi_i}
\]

Onde:
- \(y_i\) é o valor da característica para a unidade \(i\),
- \(\pi_i\) é a probabilidade de inclusão da unidade \(i\),
- \(s\) é a amostra.

Dado o plano amostral e as probabilidades de inclusão calculadas anteriormente, vamos calcular \(\hat{t}_\pi\) para cada amostra possível:

Para a amostra \(s_1 = \{1, 2\}\):
\[
\hat{t}_\pi(s_1) = \frac{y_1}{\pi_1} + \frac{y_2}{\pi_2} = \frac{2}{\frac{3}{4}} + \frac{4}{\frac{3}{4}} = \frac{2 \times 4}{3} + \frac{4 \times 4}{3} = \frac{8}{3} + \frac{16}{3} = 8
\]

Para a amostra \(s_2 = \{1, 3\}\):
\[
\hat{t}_\pi(s_2) = \frac{y_1}{\pi_1} + \frac{y_3}{\pi_3} = \frac{2}{\frac{3}{4}} + \frac{6}{\frac{1}{2}} = \frac{8}{3} + 12 = \frac{8}{3} + \frac{36}{3} = \frac{44}{3} \approx 14.67
\]

Para a amostra \(s_3 = \{2, 3\}\):
\[
\hat{t}_\pi(s_3) = \frac{y_2}{\pi_2} + \frac{y_3}{\pi_3} = \frac{4}{\frac{3}{4}} + \frac{6}{\frac{1}{2}} = \frac{16}{3} + 12 = \frac{16}{3} + \frac{36}{3} = \frac{52}{3} \approx 17.33
\]

Agora, vamos determinar a distribuição de probabilidades do estimador \(\hat{t}_\pi\):

\[
P\left(\hat{t}_\pi = 8\right) = P(s_1) = \frac{1}{2}
\]
\[
P\left(\hat{t}_\pi = \frac{44}{3} \approx 14.67\right) = P(s_2) = \frac{1}{4}
\]
\[
P\left(\hat{t}_\pi = \frac{52}{3} \approx 17.33\right) = P(s_3) = \frac{1}{4}
\]

A distribuição de probabilidades do estimador \(\hat{t}_\pi\) é dada por:

\[
\hat{t}_\pi =
\begin{cases}
8, & \text{com probabilidade } \frac{1}{2}, \\
\frac{44}{3} \approx 14.67, & \text{com probabilidade } \frac{1}{4}, \\
\frac{52}{3} \approx 17.33, & \text{com probabilidade } \frac{1}{4}.
\end{cases}
\]
}
\item[c)] Determine a variância do estimador de Horvitz-Thompson $\hat{t}_\pi$ para o total populacional da renda bruta (mensal em salários mínimos) familiar sob este plano amostral.
\end{description}
\medskip
{\scriptsize

\textbf{Res:  1c} Podemos dividir o cálculo da variância em partes menores para facilitar. Cálculo da Variância do Estimador de Horvitz-Thompson \(\hat{t}_\pi\) A variância do estimador de Horvitz-Thompson é dada por:
\[
\text{Var}(\hat{t}_\pi) = \sum_{i \in \mathcal{U}} \sum_{j \in \mathcal{U}} \left(\frac{y_i}{\pi_i}\right) \left(\frac{y_j}{\pi_j}\right) (\pi_{ij} - \pi_i \pi_j)
\]

Passo 1: Calcular os termos diagonais

Primeiro, calculamos os termos em que \(i = j\), ou seja, onde a soma envolve as mesmas unidades:
\[
\text{Var}(\hat{t}_\pi)_{\text{diagonal}} = \sum_{i \in \mathcal{U}} \left(\frac{y_i}{\pi_i}\right)^2 (1 - \pi_i)
\]

Os valores que temos são:

- \(\pi_1 = \frac{3}{4}\), \(\pi_2 = \frac{3}{4}\), \(\pi_3 = \frac{1}{2}\)

- \(y_1 = 2\), \(y_2 = 4\), \(y_3 = 6\)

Substituindo:
\[
\text{Var}(\hat{t}_\pi)_{\text{diagonal}} = \left(\frac{2}{\frac{3}{4}}\right)^2 \left(\frac{1}{4}\right) + \left(\frac{4}{\frac{3}{4}}\right)^2 \left(\frac{1}{4}\right) + \left(\frac{6}{\frac{1}{2}}\right)^2 \left(\frac{1}{2}\right)
\]

Agora, vamos calcular esses três termos separadamente:

1. Para \(i = 1\):
\[
\left(\frac{2}{\frac{3}{4}}\right)^2 \times \frac{1}{4} = \left(\frac{8}{3}\right)^2 \times \frac{1}{4} = \frac{64}{9} \times \frac{1}{4} = \frac{16}{9} \approx 1.78
\]

2. Para \(i = 2\):
\[
\left(\frac{4}{\frac{3}{4}}\right)^2 \times \frac{1}{4} = \left(\frac{16}{3}\right)^2 \times \frac{1}{4} = \frac{256}{9} \times \frac{1}{4} = \frac{64}{9} \approx 7.11
\]

3. Para \(i = 3\):
\[
\left(\frac{6}{\frac{1}{2}}\right)^2 \times \frac{1}{2} = 12^2 \times \frac{1}{2} = 144 \times \frac{1}{2} = 72
\]

Somando esses termos:
\[
\text{Var}(\hat{t}_\pi)_{\text{diagonal}} \approx 1.78 + 7.11 + 72 = 80.89
\]
Este é o valor parcial para a variância considerando apenas os termos diagonais. Vamos agora calcular os termos cruzados da variância do estimador de Horvitz-Thompson, onde \(i \neq j\).

\[
\text{Var}(\hat{t}_\pi)_{\text{cruzados}} = \sum_{i \neq j} \left(\frac{y_i}{\pi_i}\right) \left(\frac{y_j}{\pi_j}\right) (\pi_{ij} - \pi_i \pi_j)
\]
Aqui, precisamos calcular cada combinação de \(i\) e \(j\) onde \(i \neq j\). Lembrando que as probabilidades de segunda ordem \(\pi_{ij}\) foram calculadas anteriormente: \(\pi_{12} = \frac{1}{2}\), \(\pi_{13} = \frac{1}{4}\) e \(\pi_{23} = \frac{1}{4}\). O Termo para \(i = 1\) e \(j = 2\):
\[
\left(\frac{y_1}{\pi_1}\right) \left(\frac{y_2}{\pi_2}\right) (\pi_{12} - \pi_1 \pi_2) = \left(\frac{2}{\frac{3}{4}}\right) \left(\frac{4}{\frac{3}{4}}\right) \left(\frac{1}{2} - \frac{3}{4} \times \frac{3}{4}\right)
\]
\[
= \frac{8}{3} \times \frac{16}{3} \times \left(\frac{1}{2} - \frac{9}{16}\right) = \frac{128}{9} \times \left(\frac{8}{16}\right) = \frac{128}{9} \times \frac{1}{2} = \frac{64}{9} \approx 7.11
\]

O termo para \(i = 1\) e \(j = 3\): 
\[
\left(\frac{y_1}{\pi_1}\right) \left(\frac{y_3}{\pi_3}\right) (\pi_{13} - \pi_1 \pi_3) = \left(\frac{2}{\frac{3}{4}}\right) \left(\frac{6}{\frac{1}{2}}\right) \left(\frac{1}{4} - \frac{3}{4} \times \frac{1}{2}\right)
\] 
\[
= \frac{8}{3} \times 12 \times \left(\frac{1}{4} - \frac{3}{8}\right) = \frac{96}{1} \times \left(\frac{2}{8} - \frac{3}{8}\right) = \frac{96}{1} \times \left(-\frac{1}{8}\right) = -12
\]

O termo para \(i = 2\) e \(j = 3\): 
\[
\left(\frac{y_2}{\pi_2}\right) \left(\frac{y_3}{\pi_3}\right) (\pi_{23} - \pi_2 \pi_3) = \left(\frac{4}{\frac{3}{4}}\right) \left(\frac{6}{\frac{1}{2}}\right) \left(\frac{1}{4} - \frac{3}{4} \times \frac{1}{2}\right)
\] 
\[
= \frac{16}{3} \times 12 \times \left(\frac{1}{4} - \frac{3}{8}\right) = \frac{192}{1} \times \left(\frac{2}{8} - \frac{3}{8}\right) = \frac{192}{1} \times \left(-\frac{1}{8}\right) = -24
\] 
Agora, somamos os termos cruzados: 
\[
\text{Var}(\hat{t}_\pi)_{\text{cruzados}} = 7.11 - 12 - 24 = -28.89
\] 
Finalmente, somamos os termos diagonais e cruzados para obter a variância total do estimador de Horvitz-Thompson:
\[
\text{Var}(\hat{t}_\pi) \approx 80.89 - 28.89 = 52
\] 
Este é o valor da variância do estimador de Horvitz-Thompson \(\hat{t}_\pi\) para o total populacional da renda bruta sob o plano amostral fornecido.

}

\medskip 
\question~~({\bf Plano amostral AASs}) Uma amostra aleatória simples e sem substituição de 5 pessoas foi selecionada de uma população de 100 trabalhadores da empresa LInCaTech. Foram coletadas a informações sobre a Renda mensal em miles de reais (Renda) e o sexo do trabalhador. Com as informações da tabela \ref{ID:Renda} estime:

\begin{description}
    \item[a)] A renda média dos trabalhadores. Estabeleça um intervalo de 95\% para a renda média.
\item[b)] A renda total dos trabalhadores. Estabeleça um intervalo de 95\% para a renda total. 
\end{description}
{\small \begin{table}[H]
\centering
\begin{tabular}{rlr}
  \hline
ID & Sexo & Renda  \\ 
  \hline
    1 & Fem &  1 \\
    2 & Mas &  2\\
    3 & Fem &  3\\
    4 & Fem &  4\\
    5 & Mas &  5\\
   \hline
\end{tabular}
\caption{Tabela de Informações dos empregados na amostra} 
\label{ID:Renda}
\end{table}
}


{\scriptsize
	
	Os dados iniciais são o tamanho da população, $N=100$, o tamanho da amostra, $n=5$, e os valores de renda (em milhares de reais) coletados na amostra: $y_s = \{1, 2, 3, 4, 5\}$. Calculamos as estatísticas amostrais. A média amostral da renda é:
$ \bar{y} = \frac{1}{n} \sum_{i \in s} y_i = \frac{1+2+3+4+5}{5} = \frac{15}{5} = 3 $ 
A variância amostral, que mede a dispersão dos dados na amostra, é calculada como:
$ s^2 = \frac{1}{n-1} \sum_{i \in s} (y_i - \bar{y})^2 = \frac{(1-3)^2 + (2-3)^2 + (3-3)^2 + (4-3)^2 + (5-3)^2}{4} = \frac{4+1+0+1+4}{4} = \frac{10}{4} = 2.5 $

\textbf{Res:  2a}  Renda Média dos Trabalhadores

A estimativa pontual para a renda média populacional, $\mu$, usando o estimador de Horvitz-Thompson sob um plano AASs, é simplesmente a média amostral. Portanto, a renda média estimada é de R\$ 3.000,00.

Para construir o intervalo de confiança de 95\%, precisamos da variância estimada do estimador da média, $\hat{V}(\hat{\mu})$. Esta incorpora o fator de correção para populações finitas, dado que a amostragem é sem reposição:
$ \hat{V}(\hat{\mu}) = \left(1 - \frac{n}{N}\right) \frac{s^2}{n} = \left(1 - \frac{5}{100}\right) \frac{2.5}{5} = (0.95)(0.5) = 0.475 $

O erro padrão (EP) é a raiz quadrada desta variância: $EP(\hat{\mu}) = \sqrt{0.475} \approx 0.6892$.
Com um nível de confiança de 95\%, o valor crítico da distribuição normal padrão é $z \approx 1.96$. O intervalo de confiança é, então:
$ IC(\mu) = \hat{\mu} \pm z \cdot EP(\hat{\mu}) = 3 \pm 1.96 \cdot 0.6892 \approx 3 \pm 1.3508 $
Isso nos leva ao intervalo [1.6492, 4.3508].

\medskip
\textbf{Res:  2b}   Renda Total dos Trabalhadores

De forma análoga, a estimativa pontual para a renda total na população, $\tau$, é obtida expandindo a média amostral para o tamanho da população:
$ \hat{\tau} = N \cdot \bar{y} = 100 \cdot 3 = 300 $
A estimativa da renda total é de R\$ 300.000,00.

A variância estimada do estimador do total, $\hat{V}(\hat{\tau})$, está diretamente relacionada à variância do estimador da média:
$ \hat{V}(\hat{\tau}) = N^2 \cdot \hat{V}(\hat{\mu}) = 100^2 \cdot 0.475 = 10000 \cdot 0.475 = 4750 $

O erro padrão correspondente é $EP(\hat{\tau}) = \sqrt{4750} \approx 68.92$. O intervalo de 95\% de confiança para o total é:
$ IC(\tau) = \hat{\tau} \pm z \cdot EP(\hat{\tau}) = 300 \pm 1.96 \cdot 68.92 \approx 300 \pm 135.083 $
O intervalo resultante é [164.917, 435.083]. 

}


\medskip 
\question~({\bf Amostragem Bernoulli}) Seja $s$ uma amostra obtida de um plano amostral de tipo Bernoulli com probabilidades de inclusão $\pi_k = \pi$ para todo $k \in U$ (população). Seja $n_s$ o tamanho de amostra da amostra $s.$ Mostre que a probabilidade condicional de se obter $s$ dado $n_s$ é a mesma que a probabilidade obtida por uma amostragem aleatória simples sem substituição de tamanho fixado $n_s$ de $N$ (Tamanho da população). 

\medskip
{\scriptsize  {\bf Res: 3}
Considere um plano amostral Bernoulli com probabilidades de inclusão \( \pi_k = \pi \) para todo \( k \in U \). Se \( s \) é uma amostra de tamanho \( n_s \), a probabilidade de observar essa amostra específica \( s \) é dada por: 
\[
P(s \mid n_s) = \pi^{n_s} (1 - \pi)^{N - n_s}.
\] Onde \( \pi^{n_s} \) é a probabilidade de incluir exatamente essas \( n_s \) unidades na amostra, e \( (1 - \pi)^{N - n_s} \) é a probabilidade de não incluir as restantes \( N - n_s \) unidades.

Por outro lado, na amostragem aleatória simples sem substituição (AASs), a probabilidade de obter uma amostra específica \( s \) de tamanho \( n_s \) é dada por:
\[
P(s) = \frac{1}{\binom{N}{n_s}},
\] 
onde \( \binom{N}{n_s} \) é o número total de combinações possíveis de tamanho \( n_s \) a partir da população de tamanho \( N \). 
Para amostras de tamanho fixo \( n_s \), a probabilidade de se obter uma amostra \( s \) é a mesma tanto para um plano amostral Bernoulli quanto para a amostragem aleatória simples sem substituição: 
\[
P(s \mid n_s) = \frac{1}{\binom{N}{n_s}} = P(s).
\] Isso ocorre porque, para amostras de tamanho fixo \( n_s \), todas as amostras possíveis têm a mesma probabilidade de serem selecionadas, seja em um plano amostral Bernoulli ou em uma amostragem aleatória simples sem substituição.
}


\medskip
\question~({\bf Amostragem Sistemática}) Suponha uma população de 7 elementos cujos valores para a característica de interesse sejam dados por
$\mathbf{Y} = \{1, 3, 5, 7, 6, 4, 2\}.$ Calcular a variância do estimador de Horvitz-Thompson para o total populacional em um plano amostral sistemático com $a = 2$ grupos. Para esse caso específico, o plano amostral sistemático é mais eficiente do que o plano amostral aleatório simples sem reposição? Explique.

\medskip
{\scriptsize  {\bf Res: 4}
	
Para a população finita de $N=7$ num plano amostral sistemático com um intervalo de amostragem $k=2$. Este procedimento gera $k=2$ amostras potenciais. A primeira, partindo do elemento de índice 1, é $S_1 = \{y_1, y_3, y_5, y_7\} = \{1, 5, 6, 2\}$. A segunda, partindo do elemento de índice 2, é $S_2 = \{y_2, y_4, y_6\} = \{3, 7, 4\}$. A probabilidade de seleção de qualquer uma dessas amostras é $P(S_j) = 1/k = 1/2$. Como cada elemento da população pertence a exatamente uma dessas amostras, a probabilidade de inclusão de primeira ordem para qualquer elemento $i$ é constante e igual a $\pi_i = 1/2$.

O estimador de Horvitz-Thompson para o total populacional, $\tau$, é dado por $\hat{\tau}_{HT} = \sum_{i \in s} y_i / \pi_i$. O verdadeiro total populacional é $\tau = 1+3+5+7+6+4+2 = 28$. Calculando o valor do estimador para cada amostra possível, obtemos $\hat{\tau}_1 = (1+5+6+2)/(1/2) = 14 \cdot 2 = 28$ e $\hat{\tau}_2 = (3+7+4)/(1/2) = 14 \cdot 2 = 28$.

A variância do estimador é calculada como a variabilidade das estimativas em torno do verdadeiro parâmetro, ponderada pela probabilidade de cada amostra:
$ V(\hat{\tau}_{SYS}) = \sum_{j=1}^{k} P(S_j) (\hat{\tau}_j - \tau)^2 = \frac{1}{2}(28 - 28)^2 + \frac{1}{2}(28 - 28)^2 = 0 $
A variância do estimador para este plano amostral é zero, indicando que, independentemente da amostra selecionada, a estimativa será perfeitamente acurada.

Para avaliar se este plano é mais eficiente que um plano amostral aleatório simples sem reposição (AASs), comparamos sua variância com a de um estimador AAS (sem reposição) para um tamanho de amostra equivalente. O tamanho esperado da amostra sistemática é $E[n] = (4 \cdot 1/2) + (3 \cdot 1/2) = 3.5$. A variância do estimador do total sob AASs é $V(\hat{\tau}_{AAS}) = N^2 (1 - n/N) S^2/n$. A variância populacional é $S^2 = \frac{1}{N-1}\sum(y_i - \mu)^2 = \frac{28}{6} = 14/3$. Assim,
$ V(\hat{\tau}_{AAS}) = 7^2 \left(1 - \frac{3.5}{7}\right) \frac{14/3}{3.5} = 49 \cdot (0.5) \cdot \frac{14}{10.5} \approx 32.67 $
Como $V(\hat{\tau}_{SYS}) = 0 < V(\hat{\tau}_{AAS}) \approx 32.67$, concluímos que o plano sistemático é  mais eficiente.

A análise de variância (ANOVA) nos permite decompor a variabilidade total da população em duas componentes: a variância entre as amostras sistemáticas e a variância dentro delas. Neste caso, como as estimativas de total para cada amostra foram idênticas ($\hat{\tau}_1 = \hat{\tau}_2$), a variância entre as amostras é nula, resultando em $V(\hat{\tau}_{SYS})=0$. Toda a variabilidade da população está, na verdade, contida dentro de cada amostra.

Assim, o grau de homogeneidade das amostras dada pela ordenação da população ($\{1, 3, 5, 7, 6, 4, 2\}$) exibe uma tendência periódica que o intervalo de amostragem $k=2$ captura perfeitamente. Isso resulta em amostras que são internamente muito heterogêneas, pois cada uma contém valores de toda a faixa da população. Consequentemente, as amostras são extremamente homogêneas entre si. Isto indicaque o coeficiente de correlação intraclasse, $\rho$, seria fortemente negativo, indicando que elementos dentro da mesma amostra são mais diferentes entre si do que o esperado ao acaso, o que é o cenário ideal para a eficiência deste tipo de amostragem.

O efeito de desenho (deff), definido como a razão entre a variância do plano atual e a de um plano AASs de mesmo tamanho:
$ \text{deff} = \frac{V(\hat{\tau}_{SYS})}{V(\hat{\tau}_{AAS})} = \frac{0}{32.67} = 0 $
Um $\text{deff} < 1$ indica maior eficiência que a AASs. Um $\text{deff}$ de zero representa a máxima eficiência teórica, significando que o desenho amostral eliminou completamente o erro de amostragem ao se alinhar perfeitamente com a estrutura latente na população.
}








\medskip 
\question~~({\bf Amostragem estratificada}) A Tabela \ref{tab:2} contem as informações do gasto mensal em serviços públicos de uma amostra aleatória estratificada de 120 famílias na cidade de Salvador a qual foi geograficamente dividida em três estratos: Norte, Centro e Sul.

\begin{table}[!htb]
\begin{tabular}{l|ccc}
             & \multicolumn{3}{c}{Estratos}   \\ \hline
Estatísticas & Norte (1)   & Centro (2) & Sul (3)      \\ \hline
$N_h$        & 4000     & 6000     & 10000    \\
$W_h$        & 0,3      & 0,2      & 0,5      \\
$n_h$        & 40       & 36       & 44       \\
$\bar{y}_h$  & 1,2      & 2,4      & 0,6      \\
$\bar{Y}_h$  & 9600     & 7200     & 6000     \\
$s^2_h$             & 0,36     & 1,21     & 0,04     \\
${\rm Var}(\bar{y}_h)$          & 0,000993 & 0,004404 & 0,000226 \\ \hline
\end{tabular}
\caption{Informações do gasto familiar mensal em serviços públicos (em salários mínimos) a partir de uma amostra aleatória simples estratificada na cidade de Salvador.}
\label{tab:2}
\end{table}

Estime o  gasto total de toda a população e estabeleça um intervalo de confiança de 95\% para o total populacional.
 

 {\scriptsize
		
		O gasto médio da população (em salários mínimos), $\bar{y}$, é dado por:\[
		\bar{y} = \sum_{h=1}^{H} W_h \bar{y}_h = 0,3 \times 1,2 + 0,2 \times 2,4 + 0,5 \times 0,6 = 1,14
		\] 
		O gasto total estimado da população (em salários mínimos), $\hat{Y}$, é:
		\[
		\hat{Y} = N \bar{y} = (4000 + 6000 + 10000) \times 1,14 = 20.000 \times 1,14 = 22.800 
		\]
		O erro padrão da média é dado pela fórmula ajustada, incluindo o fator de correção finita:
		\[
		\text{Var}(\bar{y}) = \sum_{h=1}^{H} W_h^2 \frac{s_h^2}{n_h} \left(\frac{N_h - n_h}{N_h}\right).
		\]
		Substituindo os valores dos estratos:
		\[
		\text{Var}(\bar{y}) = (0,3)^2 \times \frac{0,36}{40} \times \frac{4000 - 40}{4000} + (0,2)^2 \times \frac{1,21}{36} \times \frac{6000 - 36}{6000} 
		\]
		\[+ (0,5)^2 \times \frac{0,04}{44} \times \frac{10000 - 44}{10000}.
		\]
		\[
		\text{Var}(\bar{y}) \approx 0,002944.
		\] Logo, 
		\[
		\text{EP}(\bar{y}) = \sqrt{0,002944} = 0,0543.
		\]
		Para construir o intervalo de confiança usando a distribuição $t$-Student, devemos calcular os graus de liberdade aproximados. Como o plano é estratificado, podemos usar a fórmula aproximada para graus de liberdade em amostragem estratificada:
		\[
		df \approx \frac{\left(\sum_{h=1}^{H} W_h^2 \frac{s_h^2}{n_h} \right)^2}{\sum_{h=1}^{H} \left(\frac{W_h^4 s_h^4}{n_h^2 (n_h - 1)}\right)}.
		\]
		Calculando este valor para os dados fornecidos, obtemos:
		\[
		df \approx 78.
		\]
		O valor crítico da distribuição t para $df = 78$ e um nível de confiança de 95\% é aproximadamente $t_{0,025} \approx 1,99$.
		O intervalo de confiança de 95\% para o gasto médio da população é dado por:
		\[
		\bar{y} \pm t_{0,025} \times \text{EP}(\bar{y}) = 1,14 \pm 1,99 \times 0,0543 = [1,032, 1,248] \text{ salários mínimos.}
		\]
		Para o gasto total da população, o intervalo de confiança de 95\% é:
		\[
		\hat{Y} \pm t_{0,025} \times \text{EP}(\hat{Y}) = 22.800 \pm 1,99 \times (20.000 \times 0,0543) = [21.724, 23.876] 
		\]
		
		Se usamos uma aproximação apenas normal, supondo que o $n$ é suficientemente grande, temos que o intervalo de confiança de 95\% para o gasto médio é:
		
		\[
		\bar{y} \pm z_{0,025} \times \text{EP}(\bar{y}) = 1,14 \pm 1,96 \times 0,0543 = [1,033, 1,264] 
		\]
		
		Para o gasto total da população, o intervalo de confiança de 95\% é:
		\[
		\hat{Y} \pm z_{0,025} \times \text{EP}(\hat{Y}) = 22.800 \pm 1,96 \times (20.000 \times 0,0543) = [20.671, 24.928] 
		\]
}


% \vspace{5cm}

\bigskip


\centering\color{cor}{\scshape boa prova}


\end{document}

% Prova
